\clearpage\fancyhead[L]{\small JUNO.}\hypertarget{native-380-lower}{}\label{native-380-lower}
\begin{center}\textbf{JUNO.}\par\vspace{4pt}\hrule\vspace{4pt}\fontsize{10}{12}\selectfont Göttingische gelehrte Anzeigen. Stück 19. Seite 189. 31. Januar 1814.\par\vspace{4pt}\hrule\vspace{4pt}\end{center}

\par\indent Als Nachtrag zu unserer kürzlich erschienenen Mitteilung über Juno (siehe Stück 204 vom vergangenen Jahr) fügen wir noch drei Beobachtungen dieses Planeten von Professor \textsc{Bessel} hinzu:

\begin{center}\fontsize{10}{12}\selectfont\setlength{\tabcolsep}{3pt}\renewcommand{\arraystretch}{1.15}\begin{tabular}{l|c|rrr|c}
1812 & \shortstack{Mittlere Zeit\\in Königsberg} & \multicolumn{3}{c|}{\shortstack{Rektaszension}} & \shortstack{Südliche\\Deklination}\\\hline
Nov. 18 & $12^{\mathrm u}\;12^{\mathrm m}\;26^{\mathrm s}5$ & $60^\circ$ & $41'$ & $21''5$ & ---\quad---\\
21 & $11\;58\;20.4$ & 60 & 6 & 42.5 & ---\quad---\\
Dez. 11 & $10\;25\;54.9$ & 56 & 39 & 0.8 & $3^\circ\;53'\;27''2$\\
\end{tabular}\end{center}

\par\indent Dies sind die ersten Früchte der neuen Königsberger Sternwarte, deren Bau vor drei Jahren begann und die Herr \textsc{Bessel} schon im vergangenen Herbst bezogen hat. Die Astronomie schuldet der Regierung den wärmsten Dank dafür, dass sie dieses für die Wissenschaft so wichtige Institut gegründet, mit vorzüglichen Instrumenten ausgestattet und selbst in diesen bedrängten Jahren so rasch vollendet hat. Jeder Monat, den ein Astronom wie \textsc{Bessel} länger auf die Möglichkeit hätte warten müssen, seine Arbeit voll zu entfalten, wäre in der Tat ein Verlust für die Wissenschaft gewesen.
